Este trabalho teve como foco o desenvolvimento de uma solução tecnológica como auxílio ao pré-diagnóstico de indícios de TDA em crianças entre 10 e 12 anos, por meio da integração entre rastreamento ocular, métricas do Teste de Desempenho Contínuo, \textit{IoT} e IA.

Os resultados obtidos demonstram que a solução funciona na etapa atual do projeto. O jogo foi implementado, integrado aos módulos de coleta de dados oculares em tempo real e capaz de gerar retorno inicial sobre o desempenho atencional dos participantes a partir do processamento automático dos dados. 

No módulo de rastreamento ocular, observou-se precisão satisfatória para a fase atual, confirmando a viabilidade técnica da abordagem. Em relação às métricas TDC, a aplicação foi adequada durante os testes realizados. Por fim, no componente de \textit{IoT}, o funcionamento foi adequado e promissor durante os testes. Entretanto, permanecem pontos de melhoria já identificados: aumento da robustez e da estabilidade da predição ocular em diferentes condições de uso, por meio de estudos de bibliotecas e técnicas de \textit{eye tracking}; aperfeiçoamento ergonômico do \textit{case} físico do dispositivo; e implementação do módulo de IA, cuja consolidação depende da confiabilidade dos dados de rastreamento ocular de entrada.

Como continuidade do estudo, devem ser realizados testes em ambiente controlado com profissionais da saúde e crianças da faixa etária-alvo, além da validação progressiva dos modelos de IA para identificação de padrões e emissão de relatórios detalhados. Conclui-se, portanto, que a proposta desenvolvida apresenta aplicabilidade prática e potencial de contribuição para auxiliar na identificação de indícios de TDA, ao mesmo tempo em que requer aprimoramentos incrementais para atingir maior maturidade técnica e operacional.